\documentclass[12pt,a4paper]{article}
\usepackage[utf8]{inputenc}
\usepackage[spanish]{babel}
\usepackage{geometry}
\usepackage{listings}
\usepackage{xcolor}
\usepackage{parskip}
\usepackage{amsmath}
\usepackage{hyperref}

\geometry{margin=2.5cm}

\lstset{
  basicstyle=\ttfamily\small,
  keywordstyle=\color{blue}\bfseries,
  commentstyle=\color{gray},
  breaklines=true, frame=single, language=Java,
  numbers=left, numberstyle=\tiny\color{gray}
}

\title{\textbf{8-Puzzle Interactivo con Heurísticas}\\
  \large Documentación Técnica}
\author{Inteligencia Artificial}
\date{}

\begin{document}
\maketitle
\tableofcontents
\newpage

\section{Descripción del Problema}
El \textbf{8-Puzzle} consiste en un tablero de 3×3 con 8 fichas numeradas del 1 al 8 y un espacio vacío. El objetivo es alcanzar el estado meta deslizando fichas al espacio vacío.

\[
\text{Estado meta:}\quad
\begin{bmatrix} 1 & 2 & 3 \\ 4 & 5 & 6 \\ 7 & 8 & \_ \end{bmatrix}
\]

\section{Representación del Estado}

El tablero se representa como un arreglo unidimensional de 9 enteros, donde el valor \texttt{0} representa el espacio vacío. Los movimientos legales desplazan el espacio vacío en una de las 4 direcciones (arriba, abajo, izquierda, derecha).

\section{Heurísticas Implementadas}

\subsection{Distancia Manhattan}
Suma de distancias horizontal y vertical de cada ficha hasta su posición meta:
\[
h_M(n) = \sum_{i=1}^{8} \left(|fila_i - fila_i^*| + |col_i - col_i^*|\right)
\]

Esta heurística es \textbf{admisible} (nunca sobreestima) y \textbf{consistente}, lo que garantiza que la búsqueda encuentre la solución óptima.

\subsection{Piezas Mal Colocadas}
Cuenta el número de fichas que no están en su posición correcta:
\[
h_P(n) = |\{i : tablero[i] \neq meta[i],\ tablero[i] \neq 0\}|
\]

Es menos informada que Manhattan, por lo que genera más nodos antes de encontrar la solución.

\section{Algoritmo de Búsqueda Heurística}

Se implementa una variante de \textbf{Best-First Search} usando una cola de prioridad ordenada por el valor heurístico $h(n)$:

\begin{lstlisting}
Nodo_BFS resolver(int heuristica) {
    PriorityQueue<Nodo_BFS> cola =
        new PriorityQueue<>(Comparator.comparingInt(
            n -> calcularH(n.getEstado(), heuristica)));
    cola.add(estadoInicial);
    visitados.add(estadoInicial);

    while (!cola.isEmpty()) {
        Nodo_BFS actual = cola.poll();
        if (esMeta(actual)) return actual;

        for (Nodo_BFS hijo : generarHijos(actual)) {
            if (!estaVisitado(hijo)) {
                visitados.add(hijo);
                cola.add(hijo);
            }
        }
    }
    return null;
}
\end{lstlisting}

\section{Cálculo de la Heurística Manhattan}

\begin{lstlisting}
int manhattan(int[] tablero) {
    int total = 0;
    for (int i = 0; i < 9; i++) {
        if (tablero[i] == 0) continue;
        int valor = tablero[i] - 1;
        int filaActual = i / 3;
        int colActual  = i % 3;
        int filaMeta   = valor / 3;
        int colMeta    = valor % 3;
        total += Math.abs(filaActual - filaMeta)
               + Math.abs(colActual  - colMeta);
    }
    return total;
}
\end{lstlisting}

\section{Generación de Movimientos}

Los movimientos válidos se determinan por la posición del espacio vacío:

\begin{lstlisting}
List<int[]> generarHijos(int[] estado) {
    int pos = posicionVacio(estado); // índice del 0
    int fila = pos / 3, col = pos % 3;
    int[][] dirs = {{-1,0},{1,0},{0,-1},{0,1}};
    List<int[]> hijos = new ArrayList<>();

    for (int[] d : dirs) {
        int nf = fila + d[0], nc = col + d[1];
        if (nf >= 0 && nf < 3 && nc >= 0 && nc < 3) {
            int[] hijo = estado.clone();
            int npos = nf*3 + nc;
            hijo[pos] = hijo[npos];
            hijo[npos] = 0;
            hijos.add(hijo);
        }
    }
    return hijos;
}
\end{lstlisting}

\section{Interfaz Gráfica}

\begin{itemize}
  \item El tablero se dibuja con \texttt{JButton[9]} en un \texttt{GridLayout(3,3)}.
  \item El botón \textbf{Mezclar} aplica 50 movimientos aleatorios válidos.
  \item Al presionar \textbf{Resolver}, se ejecuta el algoritmo y se anima la solución paso a paso.
  \item El contador muestra los pasos de la solución encontrada.
\end{itemize}

\end{document}
